\documentclass[11pt]{article}

\usepackage[margin=1.08in]{geometry}
\usepackage[T1]{fontenc}
\usepackage{lmodern}
\usepackage{amsmath,amssymb,amsthm,mathtools}
\usepackage{enumitem}
\usepackage{microtype}
\usepackage{mathrsfs}
\usepackage{hyperref}

\hypersetup{
  colorlinks=true,
  linkcolor=blue,
  citecolor=blue,
  urlcolor=blue,
  hypertexnames=false,
  pdftitle={Spectral Poisson--Dirichlet Universality for Luce Permutations under Arbitrary Rate Arrangements},
  pdfauthor={Christopher D. Long}
}

\allowdisplaybreaks

\newtheorem{theorem}{Theorem}[section]
\newtheorem{proposition}[theorem]{Proposition}
\newtheorem{lemma}[theorem]{Lemma}
\newtheorem{corollary}[theorem]{Corollary}
\newtheorem{conjecture}[theorem]{Conjecture}
\newtheorem{question}[theorem]{Question}
\theoremstyle{definition}
\newtheorem{definition}[theorem]{Definition}
\theoremstyle{remark}
\newtheorem{remark}[theorem]{Remark}

\newcommand{\Pp}{\mathbb P}
\newcommand{\Ee}{\mathbb E}
\newcommand{\Var}{\operatorname{Var}}
\newcommand{\Cov}{\operatorname{Cov}}
\newcommand{\Corr}{\operatorname{Corr}}
\newcommand{\Exp}{\operatorname{Exp}}
\newcommand{\Unif}{\operatorname{Unif}}
\newcommand{\PD}{\operatorname{PD}}
\newcommand{\GEM}{\operatorname{GEM}}
\newcommand{\one}{\mathbf 1}
\newcommand{\dd}{\,d}
\newcommand{\eps}{\varepsilon}
\newcommand{\cS}{\mathcal S}
\newcommand{\cH}{\mathcal H}
\newcommand{\cG}{\mathcal G}
\newcommand{\cJ}{\mathcal J}
\newcommand{\cC}{\mathcal C}
\newcommand{\cR}{\mathcal R}
\newcommand{\norm}[1]{\left\lVert#1\right\rVert}
\newcommand{\abs}[1]{\left\lvert#1\right\rvert}

\title{Spectral Poisson--Dirichlet Universality for Luce Permutations\\
       under Arbitrary Rate Arrangements}
\author{Christopher D. Long\thanks{Email: \texttt{galizur@gmail.com}}}
\date{20 October 2026}

\begin{document}
\maketitle

\begin{abstract}
A Luce permutation is obtained by ranking independent exponential clocks with
prescribed positive rates.  We prove an arrangement-free criterion for its
macroscopic cycles: if the empirical normalized rate law converges to a law
whose exponential observation channel has a strict mean-zero $L^2$ spectral
gap, then the size-biased cycle lengths converge in $\ell^1$ to $\GEM(1)$ and
the ranked cycle lengths converge in $\ell^1$ to $\PD(1)$.  Only the
asymptotic rate spectrum matters; the rates may be assigned to the labels in
an arbitrary deterministic order.

The proof identifies two mechanisms.  A factorization through the marginal
rate channel shows that label--rate coupling cannot worsen its spectral gap,
and an exact finite-state Markov contraction controls all late clock rows by
their total mass, without a Hilbert--Schmidt assumption.  A dyadic
Poincar\'e inequality combined with Karamata theory establishes the channel
gap for every rate law with a regularly varying lower tail of positive index.
Consequently the theorem covers all uniformly comparable arrays, every
rearrangement of the Sukhatme rates $1,2,\ldots,n$, every rearrangement of
$1^\gamma,2^\gamma,\ldots,n^\gamma$ for $\gamma>0$, and general positive
regularly varying spectra.  In particular, it resolves the macroscopic- and
longest-cycle problem for the Sukhatme model.  We also give a subsequential
criterion and a quenched extension to random rate environments.  Two further
spectral results identify a broader sufficient condition and a boundary of the
criterion: successive dyadic shell-mass ratios bounded away from zero and one
imply a strict channel gap, whereas a slowly varying lower tail has channel
norm one.
\end{abstract}

\medskip
\noindent\textbf{Keywords.}
Luce permutations; Poisson--Dirichlet distributions; exponential races;
spectral gaps; regular variation.

\section{Introduction}
\label{sec:introduction}

Let $X_{n,1},\ldots,X_{n,n}$ be independent exponential random variables with
positive rates $\theta_{n,1},\ldots,\theta_{n,n}$.  Their increasing order
defines the label-to-position permutation
\begin{equation}
  \sigma_n(i)=1+\#\{j:X_{n,j}<X_{n,i}\}.
  \label{eq:luce-permutation}
\end{equation}
This is the exponential-race representation of the Plackett--Luce ranking
model; see Luce~\cite{Luce1959}, Plackett~\cite{Plackett1975}, and
Chatterjee--Diaconis--Kim~\cite{CDK2024}.  Multiplying all rates by one
positive constant leaves the permutation unchanged.

The cycle structure of a uniform random permutation is governed by the Ewens
parameter-one law~\cite{Ewens1972}.  Its normalized ranked cycles converge to
Kingman's $\PD(1)$ distribution, while their size-biased order converges to the
$\GEM(1)$ residual-allocation law; see Kingman~\cite{Kingman1975} and
Pitman~\cite[Chapter~3]{Pitman2006}.  For Luce permutations the entry law,
permuton, local patterns, and short-cycle behavior may be highly nonuniform.
The question addressed here is whether the macroscopic cycle partition
nevertheless retains the uniform-permutation limit.

\subsection*{Recent Luce models and the cycle questions}

Chatterjee, Diaconis, and Kim~\cite{CDK2024} developed the enumerative theory
of the Tsetlin library and explicitly asked for the number of cycles, the
longest cycle, the order, and the longest increasing subsequence of a
Luce-distributed permutation.  They also singled out the decreasing weights
$n,n-1,\ldots,1$, arising from exponential spacings, and called them the
Sukhatme weights.  Borga, Chatterjee, and Diaconis~\cite{BCD2026} subsequently
proved permuton and local-limit theorems, limiting pattern laws, a quenched
Benjamini--Schramm limit, and an inversion central limit theorem.  Their
Sukhatme permuton is singular at an endpoint and its pattern statistics are
far from uniform, making the macroscopic-cycle question especially natural.

A related but different source of Poisson--Dirichlet limits is the model of
random standardized permutations studied by Guerder~\cite{Guerder2026}.  In
that model one standardizes an i.i.d. word, and a uniform upper bound on the
largest letter probability yields a $\PD(1)$ limit for large cycles.  The
present model instead starts from independent continuous clocks with a
deterministic triangular array of rates.  Our criterion is spectral, permits
rates arbitrarily close to zero, and is invariant under every deterministic
assignment of a prescribed asymptotic rate spectrum to the labels.

This arrangement invariance is a central point.  Permuton and local-statistic
limits naturally retain the coupling between labels and rates
\cite{BCD2026}; the macroscopic cycle partition proved below depends only on
the limiting marginal distribution of the normalized rates.

This paper extends the author's earlier comparable-rate theorem
\cite{LongPD2026}, whose full scope, including arrays without a limiting
empirical rate law, is recovered in Corollary~\ref{cor:comparable}.  We retain
its clock-slab and residual permanental framework, adaptive relative-prefix
estimates, rooted exploration, and deletion-regeneration argument, reproducing
the necessary proofs here.  The new ingredients are domination by the
rate-observation channel, a late-row contraction estimate that avoids
Hilbert--Schmidt assumptions, and spectral-gap criteria for endpoint-singular
rate laws.  These extend the macroscopic cycle limits to Sukhatme and
positive-index regularly varying spectra under arbitrary deterministic
arrangements.

\subsection*{The exponential channel}

Normalize
\[
  r_{n,i}=\frac{\theta_{n,i}}{\max_j\theta_{n,j}}\in(0,1]
\]
and let
\begin{equation}
  \nu_n=\frac1n\sum_{i=1}^n\delta_{r_{n,i}}
  \label{eq:empirical-rate-law}
\end{equation}
be the empirical normalized rate law.  Suppose first that
$\nu_n\Rightarrow\nu$ on $[0,1]$.  The limiting exponential channel is obtained
by sampling $R\sim\nu$ and then sampling $T\mid R=r\sim\Exp(r)$.  Its time
marginal has density
\begin{equation}
  B_\nu(t)=\int_{(0,1]} r e^{-rt}\,\nu(\dd r),
  \qquad
  \mu_\nu(\dd t)=B_\nu(t)\dd t.
  \label{eq:Bnu}
\end{equation}
When $\nu(\{0\})=0$, $\mu_\nu$ is a probability measure.  Define
\begin{equation}
  (C_\nu f)(r)=\int_0^\infty f(t)r e^{-rt}\dd t,
  \qquad
  C_\nu:L^2(\mu_\nu)\longrightarrow L^2(\nu).
  \label{eq:channel}
\end{equation}
This is conditional expectation of $f(T)$ given $R=r$.  It preserves constants
and is an $L^2$ contraction.  Its mean-zero norm
\begin{equation}
  q(\nu)=
  \norm{C_\nu}_{L^2_0(\mu_\nu)\to L^2_0(\nu)}
  \label{eq:qnu}
\end{equation}
is the maximal $L^2$ correlation between the latent rate and its observed
clock time.  It is the only global mixing parameter needed below.

Let $U_1,U_2,\ldots$ be independent $\Unif(0,1)$ variables.  We write
$\GEM(1)$ for
\[
 U_1,\quad (1-U_1)U_2,\quad
 (1-U_1)(1-U_2)U_3,\ldots,
\]
and $\PD(1)$ for its decreasing rearrangement.  Put
\[
 \cS^\downarrow
 =\left\{x_1\ge x_2\ge\cdots\ge0:
          \sum_{j\ge1}x_j=1\right\},
\]
with the $\ell^1$ metric.  If
$L_{n,1}\ge L_{n,2}\ge\cdots$ are the cycle lengths of $\sigma_n$, let
$\widetilde L_{n,1},\widetilde L_{n,2},\ldots$ be their size-biased deletion
order: at each stage choose a vertex uniformly from the residual set, record
its cycle length, delete that cycle, and repeat.

\begin{theorem}[Spectral Poisson--Dirichlet criterion]
\label{thm:main-spectral}
Assume that
\[
  \nu_n\Rightarrow\nu,
  \qquad \nu(\{0\})=0,
  \qquad q(\nu)<1.
\]
Then
\begin{equation}
 \left(\frac{\widetilde L_{n,1}}n,
       \frac{\widetilde L_{n,2}}n,\ldots\right)
 \Rightarrow \GEM(1)
 \quad\text{in }\ell^1,
 \label{eq:gem-main}
\end{equation}
and
\begin{equation}
 \left(\frac{L_{n,1}}n,
       \frac{L_{n,2}}n,\ldots\right)
 \Rightarrow \PD(1)
 \quad\text{in }\cS^\downarrow.
 \label{eq:pd-main}
\end{equation}
If $I_n$ is a uniformly chosen label and $\mathscr L_n(I_n)$ is the length of
its cycle, then
\begin{equation}
  \frac{\mathscr L_n(I_n)}n\Rightarrow\Unif(0,1).
  \label{eq:uniform-root-main}
\end{equation}
The conclusion depends only on the limiting rate marginal $\nu$; no
regularity assumption is imposed on the assignment of rates to labels.
\end{theorem}

The convergence of the empirical spectrum can be weakened to a condition on
all of its limit points.

\begin{corollary}[Subsequential spectral criterion]
\label{cor:subsequential-spectral}
Suppose that every weak subsequential limit $\nu$ of $(\nu_n)$ satisfies
$\nu(\{0\})=0$ and $q(\nu)<1$.  Then
\eqref{eq:gem-main}--\eqref{eq:uniform-root-main} hold.
\end{corollary}

\begin{proof}
Every subsequence has a further subsequence along which $\nu_n$ converges.
Theorem~\ref{thm:main-spectral} identifies the same universal limit on every
such refinement.  The subsequence principle gives the result.
\end{proof}

The criterion also has a quenched random-environment form.

\begin{corollary}[Random rate environments]
\label{cor:random-environment}
Let the positive rate vectors
$\boldsymbol\theta_n=(\theta_{n,1},\ldots,\theta_{n,n})$ be random and
independent of the exponential variables used to generate the ranking.
Suppose their empirical normalized laws $\nu_n$ converge in probability, in
the weak topology, to a deterministic law $\nu$ satisfying
$\nu(\{0\})=0$ and $q(\nu)<1$.  Conditional on
$\boldsymbol\theta_n$, the laws in \eqref{eq:gem-main} and
\eqref{eq:pd-main} converge in probability, in bounded-Lipschitz distance, to
$\GEM(1)$ and $\PD(1)$, respectively.  In particular, the corresponding
annealed convergences hold.
\end{corollary}

\begin{proof}
From every subsequence one may extract a further subsequence along which
$\nu_n\Rightarrow\nu$ almost surely.  For almost every realization of the rate
environment, Theorem~\ref{thm:main-spectral} applies to the resulting
deterministic triangular array, so the conditional bounded-Lipschitz distance
tends to zero almost surely.  The subsequence criterion for convergence in
probability proves the quenched assertion, and bounded convergence gives the
annealed one.
\end{proof}

The theorem gives the longest-cycle law explicitly.

\begin{corollary}[Longest cycle and moment convergence]
\label{cor:longest-cycle}
Under the hypotheses of Theorem~\ref{thm:main-spectral}, let $P_1$ denote the
largest part of a $\PD(1)$ partition.  Then
\begin{equation}
  \frac{L_{n,1}}n\Rightarrow P_1,
  \qquad
  \Ee\left(\frac{L_{n,1}}n\right)^s\longrightarrow\Ee P_1^s
  \quad(s>0).
  \label{eq:longest-cycle}
\end{equation}
In particular,
\[
  \frac{\Ee L_{n,1}}n\longrightarrow\lambda_{\rm GD},
\]
where $\lambda_{\rm GD}=\Ee P_1$ is the Golomb--Dickman constant.  These
conclusions answer the longest-cycle question of
Chatterjee--Diaconis--Kim~\cite{CDK2024} for the Sukhatme weights, which they
singled out as a natural example, and for every deterministic rearrangement
of $1,\ldots,n$.
\end{corollary}

\begin{proof}
The largest-coordinate map is continuous on $\cS^\downarrow$ in the
$\ell^1$ metric.  Equation~\eqref{eq:pd-main} therefore gives the distributional
limit.  Since $0\le L_{n,1}/n\le1$, every positive moment converges.  The
identification of $\Ee P_1$ with the Golomb--Dickman constant is classical;
see Shepp and Lloyd~\cite{SheppLloyd1966}.
\end{proof}

The theorem includes the uniformly comparable regime: support in $[c,1]$
gives a common exponential component and hence
$q(\nu)\le\sqrt{1-c}<1$.  The new phenomenon is that support may reach zero.
Our principal sufficient condition is regular variation.

\begin{theorem}[Regularly varying channel gap]
\label{thm:rv-gap}
Let $\nu$ be a probability law on $(0,1]$, and put
\[
 F_\nu(x)=\nu((0,x]).
\]
If $F_\nu$ is regularly varying at zero with index $\alpha>0$, then
$q(\nu)<1$.  In particular, for every $\alpha>0$,
\begin{equation}
  \nu_\alpha(\dd r)=\alpha r^{\alpha-1}\one_{(0,1)}(r)\dd r
  \label{eq:power-law-law}
\end{equation}
has a strict channel gap.
\end{theorem}

\begin{corollary}[Power-law rate spectra]
\label{cor:power-law-spectra}
Fix $\gamma>0$.  Suppose that, for each $n$, the multiset of rates is
\[
  \{1^\gamma,2^\gamma,\ldots,n^\gamma\},
\]
assigned to the labels in an arbitrary deterministic order.  Then
\eqref{eq:gem-main}--\eqref{eq:uniform-root-main} and
\eqref{eq:longest-cycle} hold.  The case $\gamma=1$ includes every
rearrangement of the Sukhatme rates.
\end{corollary}

\begin{corollary}[Regularly varying spectra]
\label{cor:regularly-varying}
Let $(a_m)_{m\ge1}$ be a positive regularly varying sequence of index
$\gamma>0$, and for each $n$ assign the multiset
$\{a_1,\ldots,a_n\}$ to the labels in an arbitrary deterministic order.
Then \eqref{eq:gem-main}--\eqref{eq:uniform-root-main} and
\eqref{eq:longest-cycle} hold.
\end{corollary}

Regular variation is used to identify the limiting empirical spectrum and to
prove its spectral gap.  We use the uniform convergence theorem, Potter bounds,
and Karamata's Abelian theorem from
Bingham--Goldie--Teugels~\cite{BGT1987}.  Standard compactness, maximal
inequalities, and weak-convergence facts are taken from
Kallenberg~\cite{Kallenberg2002}, Durrett~\cite{Durrett2019},
Billingsley~\cite{Billingsley1999}, and Dudley~\cite{Dudley2002}.

Two structural points drive the proof.  First, a marked subsequential limit
remembers the arrangement of rates along the labels, but its channel factors
through $C_\nu$ and therefore cannot have a worse mean-zero norm.  Second, late
clock rows are controlled by the exact finite-state inequality
\[
  \sum_b\frac{(\sum_{a\in D}p_aP_{ab})^2}{p_b}\le p(D)
\]
for every stochastic matrix $P$ with positive stationary law $p$.  This replaces the
global Hilbert--Schmidt estimate that would fail for the Sukhatme kernel.
Sections~\ref{sec:channel}--\ref{sec:regeneration} develop the spectral and
regeneration argument, Section~\ref{sec:rv-gap} proves the regularly varying
gap, Section~\ref{sec:applications} gives the principal applications, and
Section~\ref{sec:boundary} proves the dyadic-doubling criterion and the slowly
varying boundary theorem before recording the remaining spectral and multiscale
questions.

\section{The limiting exponential channel}
\label{sec:channel}

\subsection{Marked subsequential limits}

Fix an arbitrary subsequence.  Since $[0,1]^2$ is compact, after passage to a
further subsequence the marked empirical measures
\begin{equation}
  \eta_n=\frac1n\sum_{i=1}^n\delta_{(i/n,r_{n,i})}
  \label{eq:marked-empirical}
\end{equation}
converge weakly to a probability measure $\eta$ on $[0,1]^2$.  Its first
marginal is Lebesgue measure and its second marginal is $\nu$.  Choose a
disintegration
\begin{equation}
  \eta(\dd x,\dd r)=\dd x\,\eta_x(\dd r).
  \label{eq:disintegration}
\end{equation}
Since $\nu(\{0\})=0$, we have $\eta_x(\{0\})=0$ for almost every $x$.

Put
\begin{equation}
  A(t)=\int_{(0,1]}(1-e^{-rt})\nu(\dd r),
  \qquad B(t)=A'(t)=B_\nu(t).
  \label{eq:A-and-B}
\end{equation}
Then $A$ is continuous and strictly increasing from $0$ to $1$.  Let
$\tau=A^{-1}:(0,1)\to(0,\infty)$ and define
\begin{equation}
  g_x(t)=\int_{(0,1]}r e^{-rt}\eta_x(\dd r),
  \qquad
  \rho(x,y)=\frac{g_x(\tau(y))}{B(\tau(y))}.
  \label{eq:marked-kernel}
\end{equation}

\begin{lemma}[Bistochastic marked kernel]
\label{lem:bistochastic}
The kernel $\rho$ is bistochastic up to null sets:
\[
  \int_0^1\rho(x,y)\dd y=1
  \quad\text{for a.e. }x,
  \qquad
  \int_0^1\rho(x,y)\dd x=1
  \quad\text{for a.e. }y.
\]
\end{lemma}

\begin{proof}
Under the change of variables $y=A(t)$, $\dd y=B(t)\dd t$, and therefore
\[
 \int_0^1\rho(x,y)\dd y
 =\int_0^\infty g_x(t)\dd t
 =\int_{(0,1]}\eta_x(\dd r)=1.
\]
Also $\int_0^1g_x(t)\dd x=B(t)$, which proves the column identity.
\end{proof}

Let $K_\eta$ be the integral operator with kernel $\rho$ on $L^2[0,1]$.
The next factorization explains why the arrangement of rates cannot worsen the
spectral contraction of their marginal law.

\begin{lemma}[Channel domination]
\label{lem:channel-domination}
On mean-zero subspaces,
\begin{equation}
  \norm{K_\eta}_{L^2_0\to L^2_0}\le q(\nu).
  \label{eq:marked-gap}
\end{equation}
\end{lemma}

\begin{proof}
Let $U:L^2(0,1;\dd y)\to L^2(\mu_\nu)$ be the unitary map
$(Uf)(t)=f(A(t))$.  Define
\[
 (C_\eta h)(x)=\int_0^\infty h(t)g_x(t)\dd t
\]
after which $K_\eta=C_\eta U$.
For $k\in L^2(\nu)$ put
\[
  (Jk)(x)=\int k(r)\eta_x(\dd r).
\]
Jensen's inequality gives
\[
 \norm{Jk}_{L^2(\dd x)}^2
 \le\int_0^1\int |k(r)|^2\eta_x(\dd r)\dd x
 =\norm{k}_{L^2(\nu)}^2.
\]
Fubini's theorem shows that $C_\eta=JC_\nu$.  Both factors preserve constants
and means.  Hence the mean-zero norm of $C_\eta$ is at most $q(\nu)$, and the
unitary map $U$ does not change the norm.
\end{proof}

\begin{remark}
The quantity $q(\nu)$ is the maximal $L^2$ correlation between the latent rate
$R$ and its exponential observation $T$.  We use only the operator definition
\eqref{eq:qnu}; no compactness of the channel is assumed.
\end{remark}

\section{Clock slabs and residual permanental matchings}
\label{sec:slabs}

Work along the marked subsequence fixed in Section~\ref{sec:channel}.  Replace
the original clocks by
\[
  T_{n,i}=(\max_j\theta_{n,j})X_{n,i}\sim\Exp(r_{n,i});
\]
their order, and hence $\sigma_n$, is unchanged.

Fix $T>0$ and a finite partition
\begin{equation}
  0=t_0<t_1<\cdots<t_d=T,
  \qquad \max_a(t_a-t_{a-1})\le h,
  \label{eq:time-partition}
\end{equation}
and append $I_\star=(T,\infty)$.  Put
\[
 I_a=(t_{a-1},t_a],\qquad
 H_{n,a}=\{i:T_{n,i}\in I_a\},\qquad
 m_{n,a}=|H_{n,a}|.
\]
The tail state is defined analogously.  If
$M_{n,a}=\sum_{c\le a}m_{n,c}$, the ranks occupied by $H_{n,a}$ form the
interval
\[
 Q_{n,a}=\{M_{n,a-1}+1,\ldots,M_{n,a}\}.
\]
Define
\begin{equation}
 N_{n,ab}=|Q_{n,a}\cap H_{n,b}|,
 \qquad p_{n,a}=\frac{m_{n,a}}n,
 \qquad P_{n,ab}=\frac{N_{n,ab}}{m_{n,a}}.
 \label{eq:finite-slab-matrix}
\end{equation}
On the high-probability events used below all rows are nonempty.  The matrix
$P_n$ is stochastic and $p_nP_n=p_n$.

Let
\[
 J_a=(A(t_{a-1}),A(t_a)],\qquad J_\star=(A(T),1],
 \qquad p_a=|J_a|,
\]
and
\begin{equation}
 q_{ab}=\int_{J_a\times J_b}\rho(x,y)\dd x\dd y,
 \qquad P_{ab}=q_{ab}/p_a.
 \label{eq:limit-slab-matrix}
\end{equation}
Every $p_a$ and $q_{ab}$ is positive.

\subsection{The slab law}

For a partition endpoint $U$, let $D(U)$ denote the union of slab states whose
clock intervals lie in $(U,\infty)$.  Define the limiting late-clock mass
\begin{equation}
  \overline\nu(U)=\int_{(0,1]}e^{-rU}\nu(\dd r).
  \label{eq:tail-mass}
\end{equation}
By dominated convergence, $\overline\nu(U)\downarrow0$.

\begin{lemma}[Slab law and finite-state spectral gap]
\label{lem:slab-law}
For every fixed partition,
\[
 p_{n,a}\xrightarrow{\Pp}p_a,
 \qquad
 \frac{N_{n,ab}}n\xrightarrow{\Pp}q_{ab}
\]
simultaneously in $a,b$.  There are deterministic $\eps_n\downarrow0$ and
events $\cG_n$ with $\Pp(\cG_n)\to1$ such that, on $\cG_n$:
\begin{enumerate}[label=(\roman*)]
\item all these convergences have error at most $\eps_n$, and every
      $m_{n,a}$ and $N_{n,ab}$ is bounded below by a positive constant times
      $n$;
\item for every fixed $q_0\in(q(\nu),1)$ and all sufficiently large $n$,
\begin{equation}
 \norm{P_n}_{L^2_0(p_n)\to L^2_0(p_n)}\le q_0;
 \label{eq:finite-gap}
\end{equation}
\item if $S$ is a partition endpoint and $J_b\subset[0,A(S)]$, then
\begin{equation}
  P_{n,ab}\le 2M_Sp_{n,b}
  \quad\text{for every }a,
  \qquad M_S=\frac1{B(S)};
  \label{eq:inner-column}
\end{equation}
\item for every partition endpoint $U$,
\begin{equation}
  p_n(D(U))\le \overline\nu(U)+\eps_n.
  \label{eq:finite-tail-mass}
\end{equation}
\end{enumerate}
\end{lemma}

\begin{proof}
For a fixed time cell $I_b=(u,v]$, the indicators
$\one_{\{T_{n,i}\in I_b\}}$ are independent and have means
$e^{-r_{n,i}u}-e^{-r_{n,i}v}$.  Kolmogorov's maximal inequality gives
\[
 \max_{k\le n}\left|\sum_{i\le k}
 \left(\one_{\{T_{n,i}\in I_b\}}
       -\Ee\one_{\{T_{n,i}\in I_b\}}\right)\right|=o_{\Pp}(n).
\]
Weak convergence of the marked empirical measures gives, uniformly in
$x\in[0,1]$,
\[
 \frac1n\sum_{i\le nx}\Ee\one_{\{T_{n,i}\in I_b\}}
 \longrightarrow
 \int_0^x\int
  (e^{-ru}-e^{-rv})\eta_z(\dd r)\dd z.
\]
This proves the slab-size limit and convergence of the endpoints of
$Q_{n,a}/n$ to those of $J_a$.  Evaluating the same uniform partial-sum limit
at those random endpoints yields
\[
 \frac{N_{n,ab}}n\longrightarrow
 \int_{J_a}\int(e^{-rt_{b-1}}-e^{-rt_b})\eta_x(\dd r)\dd x=q_{ab}.
\]
The last equality follows by integrating \eqref{eq:marked-kernel} over $J_b$
and using $y=A(t)$.

The matrix $P$ is the compression of $K_\eta$ to functions constant on the
partition cells $J_a$.  Conditional expectation is a contraction, so
Lemma~\ref{lem:channel-domination} gives
$\norm{P}_{L^2_0(p)\to L^2_0(p)}\le q(\nu)$.  Entrywise convergence of the
conjugated finite matrices
$D_{p_n}^{1/2}P_nD_{p_n}^{-1/2}$ proves \eqref{eq:finite-gap}.

If $y\le A(S)$, then $\tau(y)\le S$ and
\[
 \rho(x,y)=\frac{g_x(\tau(y))}{B(\tau(y))}
 \le\frac1{B(S)}=M_S,
\]
because $0<r\le1$ and $B$ is decreasing.  Hence
$P_{ab}\le M_Sp_b$, and entrywise convergence gives
\eqref{eq:inner-column}.  Finally,
\[
 p_n(D(U))=\frac1n\sum_{i=1}^n\one_{\{T_{n,i}>U\}}
 \xrightarrow{\Pp}\int e^{-rU}\nu(\dd r)=\overline\nu(U).
\]
Since the partition is finite, all statements can be placed on one sequence of
events $\cG_n$ with deterministic errors $\eps_n\downarrow0$.
\end{proof}

\subsection{Conditional permanental matchings}

Condition on the complete slab data: the sets $H_{n,a}$, the rank intervals
$Q_{n,a}$, and the ordered but unlabeled clock values
\[
 s_{a,1}<\cdots<s_{a,m_{n,a}}
\]
in every slab, including the tail slab.  The tail row will never be assigned a
flatness estimate; conditioning on its ordered values is used only for exact
residual factorization.  Write
$Q_{n,a}=\{q_{a,1}<\cdots<q_{a,m_{n,a}}\}$ and define
$\phi_a:H_{n,a}\to[m_{n,a}]$ by
$\sigma_n(i)=q_{a,\phi_a(i)}$.  For a bijection $\phi$,
independence of the exponential densities gives
\begin{equation}
 \Pp(\phi_a=\phi\mid\text{slab data})
 =\frac{\prod_{i\in H_{n,a}}e^{-r_{n,i}s_{a,\phi(i)}}}
 {\operatorname{per}(e^{-r_{n,i}s_{a,j}})_{i,j}}.
 \label{eq:permanental-row}
\end{equation}
The row matchings are conditionally independent.

\begin{lemma}[Residual permanental factorization]
\label{lem:residual-factorization}
After conditioning on any compatible partial collection of exposed matching
edges, delete the exposed source and target vertices in each row.  The remaining
row matchings are conditionally independent, and each has the permanental law
\eqref{eq:permanental-row} on its residual submatrix.
\end{lemma}

\begin{proof}
The conditional density is a product over rows and, within a row, a product
over matching edges.  Conditioning on fixed edges removes their constant
factors and restricts the normalization to bijections of the remaining rows and
columns, whose sum is the corresponding permanent.
\end{proof}

\begin{lemma}[Hereditary flatness]
\label{lem:flatness}
Let $A=(a_{ij})$ be a positive $m\times m$ matrix satisfying
\[
 e^{-\epsilon}\le a_{ij}/a_{ik}\le e^\epsilon
\]
for every row $i$ and columns $j,k$.  Under the permanental matching law
proportional to $\prod_i a_{i,\phi(i)}$, after conditioning on an arbitrary
partial matching, every remaining edge probability lies in
\[
  [e^{-2\epsilon}/s,e^{2\epsilon}/s],
\]
where $s$ rows and columns remain.  For a clock slab of width at most $h$ this
holds with
\begin{equation}
  \beta=2h.
  \label{eq:beta}
\end{equation}
\end{lemma}

\begin{proof}
Work in the current residual submatrix.  For two remaining columns $j,k$,
\[
 \frac{\Pp(\phi(i)=j)}{\Pp(\phi(i)=k)}
 =\frac{a_{ij}\operatorname{per}A_{-i,-j}}
       {a_{ik}\operatorname{per}A_{-i,-k}}.
\]
Map each term of the first minor to a term of the second by locating the row
sent to $k$ and sending that row to $j$.  Corresponding products differ by a
factor in $[e^{-\epsilon},e^\epsilon]$; the displayed prefactor contributes the
same range.  Thus pairwise edge-probability ratios lie in
$[e^{-2\epsilon},e^{2\epsilon}]$, and the individual bounds follow because the
$s$ probabilities sum to one.  In a slab,
$e^{-h}\le e^{-r_{n,i}(s_j-s_k)}\le e^h$.
\end{proof}

\begin{lemma}[Predictable exposure]
\label{lem:predictable-exposure}
Suppose that before each new edge is exposed, its source is measurable with
respect to the slab data and all previously exposed edges.  Conditional on the
current filtration, its target has the one-edge marginal of the residual
permanental matching.  If $s$ targets remain in a finite slab, each target has
conditional probability in
\[
  [e^{-\beta}/s,e^\beta/s].
\]
If one distinguished target is forbidden, each of the other $s-1$ targets has
conditional probability in
$[e^{-\beta}/(s-1),e^\beta/(s-1)]$.
\end{lemma}

\begin{proof}
Apply Lemma~\ref{lem:residual-factorization} before the new edge is revealed and
then Lemma~\ref{lem:flatness}.  Conditioning away from a distinguished target
renormalizes the other probabilities without changing their pairwise ratios.
\end{proof}

\subsection{A universal contraction for late rows}

The following elementary inequality replaces every global Hilbert--Schmidt
estimate.

\begin{lemma}[Markov contraction of a row block]
\label{lem:row-block-contraction}
Let $P$ be a finite stochastic matrix with positive stationary probability
vector $p$.
For any set $D$ of row states, put
\[
  s_b=\sum_{a\in D}p_aP_{ab}.
\]
Then
\begin{equation}
  \sum_b\frac{s_b^2}{p_b}\le p(D).
  \label{eq:row-block-contraction}
\end{equation}
\end{lemma}

\begin{proof}
Stationarity gives $0\le s_b\le p_b$, and stochasticity gives
$\sum_b s_b=p(D)$.  Therefore
\[
 \sum_b\frac{s_b^2}{p_b}\le\sum_b s_b=p(D).
\]
\end{proof}

\begin{remark}
The estimate is exact at finite $n$, requires no reversibility, and remains
valid when the limiting kernel is not Hilbert--Schmidt.  For the ordered
Sukhatme coupling $r(x)=x$, with $\nu(\dd r)=\dd r$ on $(0,1)$, the squared
Hilbert--Schmidt integral truncated at clock time $T$ satisfies
\begin{align*}
 \int_0^1\int_0^{A(T)}\rho(x,y)^2\dd y\dd x
 &=\int_0^T
   \frac{\int_0^1r^2e^{-2rt}\dd r}{B_\nu(t)}\dd t\\
 &=\frac14\log T+O(1),\qquad T\to\infty.
\end{align*}
Indeed, $B_\nu(t)=(1-(1+t)e^{-t})/t^2$, and the integrand in the time
integral is $1/(4t)+O(e^{-t})$ as $t\to\infty$.  The decreasing ordered
coupling gives the same integral by reflection.  Thus the limiting kernel is
not Hilbert--Schmidt, while Lemma~\ref{lem:row-block-contraction} still controls
late rows by the square root of their total mass.
\end{remark}

\subsection{A relative-prefix estimate}

We use the following adaptive sampling lemma.  It is included to make the
regeneration argument self-contained.

\begin{lemma}[Relative prefix law]
\label{lem:relative-prefix}
Fix $d\ge1$ and $c_*>0$.  There are absolute constants
$\beta_0>0$, $C<\infty$, and deterministic sequences
$\chi_m\downarrow0$, $\pi_m\downarrow0$, depending only on $d,c_*$, with the
following property.  Partition $m$ targets into $d$ categories of sizes
$m_1,\ldots,m_d$, where $m_b\ge c_*m$.  Select targets sequentially without
replacement in an arbitrary predictable adaptive source order.  Suppose that,
whenever $s$ targets remain, every individual remaining target has conditional
selection probability in
\[
  [e^{-\beta}/s,e^\beta/s],
  \qquad 0\le\beta\le\beta_0.
\]
If $X_b(r)$ is the number of category-$b$ targets among the first $r$
selections, then
\begin{equation}
 \Pp\left(
 \max_{b\le d}\max_{0\le r\le m}
 \frac{|X_b(r)-rm_b/m|}{m_b}>C\beta+\chi_m
 \right)\le\pi_m.
 \label{eq:relative-prefix}
\end{equation}
The conclusion is unchanged, after harmless modification of the constants, if
a fixed number of distinguished targets is deleted or forbidden.
\end{lemma}

\begin{proof}
Fix one category and write
\[
 N_r=m_b-X_b(r),\qquad s_r=m-r.
\]
If $Y_{r+1}$ indicates that the next target lies in this category, then
\begin{equation}
 \Ee(Y_{r+1}\mid\mathscr F_r)
 \in\left[e^{-\beta}\frac{N_r}{s_r},
          e^\beta\frac{N_r}{s_r}\right].
 \label{eq:prefix-drift-prob}
\end{equation}
Set $a_m=m_b^{2/3}$ and stop when $N_r<a_m$ or $s_r<a_m$.  Before that time
put $Z_r=\log(N_r/s_r)$.  Taylor expansion gives
\[
 \left|\Ee(Z_{r+1}-Z_r\mid\mathscr F_r)\right|
 \le C\frac\beta{s_r}
    +C\left(\frac1{s_r^2}+\frac1{s_rN_r}\right),
\]
and the conditional variance of the centered increment is at most
$C/(s_rN_r)$.  Since $m_b\ge c_*m$,
\[
 \sum_{r<\tau}\frac1{s_rN_r}
 \le C_{c_*}\frac{1+\log m}{m^{2/3}}=:v_m.
\]
Doob's $L^2$ maximal inequality shows that, outside an event of probability
$O(v_m^{1/2})$, the stopped martingale part is uniformly
$O(v_m^{1/4})$.  Summing the drift and writing $u=s_r/m$ yields
\[
 e^{-\eta_m}u^{1+C\beta}
 \le\frac{N_r}{m_b}
 \le e^{\eta_m}u^{1-C\beta},
 \qquad
 \eta_m=O(v_m^{1/4}+v_m).
\]
Choose $\beta_0$ so that $C\beta_0\le1/2$.  The elementary bound
\[
 \sup_{0\le u\le1}|u^{1+a}-u|\le C|a|,
 \qquad |a|\le1/2,
\]
then gives
\[
 \left|\frac{N_r}{m_b}-\frac{s_r}{m}\right|
 \le C\beta+o(1)
\]
uniformly before the stopping time.

If $s_\tau<a_m$, the upper power bound forces $N_\tau/m_b=o(1)$.  If instead
$N_\tau<a_m$, the lower power bound forces $s_\tau/m=o(1)$.  Both normalized
remaining counts are nonincreasing after $\tau$, so the same estimate holds for
all later prefixes.  A union bound over the fixed number of categories proves
\eqref{eq:relative-prefix}.  Deleting or forbidding finitely many targets
changes all relevant counts by $O(1)$ only.
\end{proof}

In all later applications the fixed mesh is chosen so that
$\beta=2h\le\beta_0$.

\section{Scale-sensitive orbit balance}
\label{sec:orbit-balance}

A residual history $\cH$ consists of the slab data, all roots selected so far,
and the matching edges in a collection of fully exposed and deleted cycles; it
contains no additional information about unexposed row matchings.  Let
$R\subset[n]$ be the remaining vertex set, $N=|R|$, and $r=N/n$.  Put
\[
 m_a^R=|H_{n,a}\cap R|,
 \qquad
 N_{ab}^R=|Q_{n,a}\cap H_{n,b}\cap R|.
\]
For every deleted cycle $C$ and every slab $a$,
\begin{equation}
 |C\cap H_{n,a}|=|C\cap Q_{n,a}|,
 \label{eq:cycle-row-balance}
\end{equation}
because $\sigma_n$ maps the first set bijectively onto the second.  Hence
$m_a^R=|Q_{n,a}\cap R|$.

Fix a partition endpoint $U$.  A history is $(U,\delta,\xi)$-regular if
$N\ge\xi n$, $0\le\delta<1/4$, and, for every row contained in $[0,U]$ and
every target state $b$,
\begin{equation}
 \left|\frac{m_a^R}{r m_{n,a}}-1\right|\le\delta,
 \qquad
 \left|\frac{N_{ab}^R}{rN_{n,ab}}-1\right|\le\delta.
 \label{eq:regular-history}
\end{equation}
The initial history is $(U,0,1)$-regular.

An admissible exploration reveals one outgoing edge at a time, with the next
source measurable before the edge is exposed.  We allow a fixed number of
distinguished targets to be forbidden.  Lemma~\ref{lem:predictable-exposure}
applies to every actual cycle exploration and every root-avoiding exploration
below.

Consider an injective path segment $v_0,\ldots,v_\ell$, allowing
$v_\ell=v_0$ only at a terminal closing step.  Put
\[
 R_a(\ell)=\#\{0\le j<\ell:v_j\in H_{n,a}\},
\]
and
\[
 R_{ab}(\ell)=\#\{0\le j<\ell:
 v_j\in H_{n,a},\ v_{j+1}\in H_{n,b}\}.
\]
The deterministic flow identity is
\begin{equation}
 R_b(\ell)-\sum_aR_{ab}(\ell)
 =\one_{\{v_0\in H_{n,b}\}}-
  \one_{\{v_\ell\in H_{n,b}\}}.
 \label{eq:flow-identity}
\end{equation}

\begin{lemma}[Scale-sensitive orbit balance]
\label{lem:orbit-balance}
Fix partition endpoints $0<S<U\le T$, $\xi>0$, and
$0\le\delta<1/4$.  There are deterministic $\chi_n,\pi_n\downarrow0$ such
that, on $\cG_n$, uniformly over every $(U,\delta,\xi)$-regular history and
every admissible exploration, with conditional probability at least
$1-\pi_n$ the following statements hold simultaneously at every defined
prefix length $\ell\ge1$.

For every row $a$ contained in $[0,U]$ and every target state $b$,
\begin{equation}
 \left|R_{ab}(\ell)-\frac{R_a(\ell)}{m_a^R}N_{ab}^R\right|
 \le(C\beta+\chi_n)N_{ab}^R.
 \label{eq:row-prefix}
\end{equation}
Define
\begin{equation}
 \Delta_n(U,\delta,h)
 =\delta+\beta+\chi_n+\sqrt{p_n(D(U))}.
 \label{eq:Delta}
\end{equation}
Then, for every inner state $b$ with $J_b\subset[0,A(S)]$,
\begin{equation}
 \left|\frac{R_b(\ell)}{\ell p_{n,b}}-1\right|
 \le C_{S,\xi}\frac n\ell\Delta_n(U,\delta,h)
      +\frac{C_{\cJ}}\ell.
 \label{eq:scale-sensitive-balance}
\end{equation}
In particular, for every fixed $\lambda>0$, the right-hand side is
$C_{S,\xi,\lambda}\Delta_n+o(1)$ uniformly for $\ell\ge\lambda n$.
\end{lemma}

\begin{proof}
On $\cG_n$, every cell in the fixed finite partition has size bounded below by
a positive constant times $n$.  Regularity gives the same bound for the
residual cells.  Apply Lemmas~\ref{lem:predictable-exposure} and
\ref{lem:relative-prefix} row by row, extending the queried source sequence
predictably after the orbit ends if necessary.  A finite union bound gives
\eqref{eq:row-prefix} with sequences $\chi_n,\pi_n$ uniform over all stated
histories and explorations.

Suppress $\ell$ and put
\[
 P=P_n,
 \qquad q_a=R_a/\ell,
 \qquad E_{ab}=R_{ab}-R_aP_{ab}.
\]
For $a\notin D(U)$, regularity and \eqref{eq:row-prefix} give
\begin{equation}
 |E_{ab}|\le C(\delta+\beta+\chi_n)N_{n,ab}.
 \label{eq:good-row-error}
\end{equation}
For $a\in D(U)$, injectivity gives
$0\le R_{ab}\le N_{n,ab}$ and
$0\le R_aP_{ab}\le N_{n,ab}$, hence
\begin{equation}
 |E_{ab}|\le N_{n,ab}.
 \label{eq:bad-row-error}
\end{equation}
Let
\[
 G_b=\sum_{a\notin D(U)}E_{ab},
 \qquad H_b=\sum_{a\in D(U)}E_{ab}.
\]
Since $\sum_aN_{n,ab}=np_{n,b}$,
\begin{equation}
 \norm{G/(\ell p_n)}_{L^2(p_n)}
 \le C\frac n\ell(\delta+\beta+\chi_n).
 \label{eq:good-L2}
\end{equation}
For the bad rows, \eqref{eq:bad-row-error} and
Lemma~\ref{lem:row-block-contraction} give
\begin{align}
 \sum_b\frac{H_b^2}{m_{n,b}}
 &\le n\sum_b
 \frac{(\sum_{a\in D(U)}p_{n,a}P_{n,ab})^2}{p_{n,b}}
 \le n p_n(D(U)).
 \label{eq:bad-L2-raw}
\end{align}
Consequently
\begin{equation}
 \norm{H/(\ell p_n)}_{L^2(p_n)}
 \le\frac n\ell\sqrt{p_n(D(U))}.
 \label{eq:bad-L2}
\end{equation}

Let $d_b$ be the right-hand side of \eqref{eq:flow-identity}; it has at most
two nonzero coordinates of absolute value at most one.  Dividing the flow
identity by $\ell$ gives
\[
 q-qP=e,
 \qquad e_b=(d_b+G_b+H_b)/\ell.
\]
Put $u_a=q_a/p_{n,a}-1$.  Since $p_nP=p_n$ and
$\sum_ap_{n,a}u_a=0$,
\begin{equation}
 u-P^*u=e/p_n.
 \label{eq:resolvent}
\end{equation}
By \eqref{eq:finite-gap},
\begin{equation}
 \norm{u}_{L^2(p_n)}
 \le(1-q_0)^{-1}\norm{e/p_n}_{L^2(p_n)}
 \le C\frac n\ell\Delta_n+\frac{C_{\cJ}}\ell.
 \label{eq:u-L2}
\end{equation}

For an inner state $b$, \eqref{eq:inner-column} and stationarity imply
\[
 \sum_ap_{n,a}\left(\frac{P_{n,ab}}{p_{n,b}}\right)^2
 \le 2M_S.
\]
Thus $|(P^*u)_b|\le\sqrt{2M_S}\norm{u}_{L^2(p_n)}$.  Directly from
\eqref{eq:good-row-error}, \eqref{eq:bad-row-error}, and
\eqref{eq:inner-column},
\[
 \left|\frac{e_b}{p_{n,b}}\right|
 \le C_S\frac n\ell
 \bigl(\delta+\beta+\chi_n+p_n(D(U))\bigr)
 +\frac{C_{\cJ}}\ell.
\]
Now use $u_b=(e/p_n)_b+(P^*u)_b$ and
$p_n(D(U))\le\sqrt{p_n(D(U))}$ to obtain
\eqref{eq:scale-sensitive-balance}.
\end{proof}

\section{Rooted cycles and regeneration}
\label{sec:rooted}

For a vertex $v$, let $c(v)$ be the row index for which $v\in Q_{n,c(v)}$.
An orbit can hit $v$ only when its current source lies in $H_{n,c(v)}$.

\subsection{Root avoidance}

\begin{lemma}[Root-avoiding change of measure]
\label{lem:root-avoiding}
Fix a residual history $\cH$ and a distinguished remaining target $v$.
Explore the path from $v$.  At each step, use the original residual one-edge
kernel conditional on the complete exposed history.  At its $j$th opportunity
to hit $v$, indexed from zero, let $\lambda_j$ be that kernel's probability of
selecting $v$.  Define the auxiliary law $\Pp^\circ$ by conditioning each such
transition away from $v$, and stopping if no allowed target remains.

Let $\tau^\circ$ be the number of nonclosing edges completed before this
stopping time.  On $\{\tau^\circ\ge k\}$, let $S(k)$ be the number of closure
opportunities among the first $k$ completed edges, and define
\begin{equation}
 W_k=
 \begin{cases}
  \displaystyle\prod_{j=0}^{S(k)-1}(1-\lambda_j),
       &\tau^\circ\ge k,\\[1ex]
  0,   &\tau^\circ<k.
 \end{cases}
 \label{eq:root-survival-weight}
\end{equation}
Then
\begin{equation}
 \Pp(|\cC(v)|>k\mid\cH,v)
 =\Ee^\circ W_k.
 \label{eq:root-rn}
\end{equation}
If the closure row $c(v)$ is a finite clock slab of width at most $h$ and
initially has $m$ available targets, including $v$, then, while defined,
\begin{equation}
 \frac{e^{-\beta}}{m-j}\le\lambda_j\le
 \frac{e^\beta}{m-j}.
 \label{eq:root-hazard}
\end{equation}
\end{lemma}

\begin{proof}
At a closure opportunity with $\lambda_j<1$, the original transition law is a
mixture of an atom of mass $\lambda_j$ at $v$ and the avoid-$v$ law of mass
$1-\lambda_j$.  Multiplying these sequential Radon--Nikodym factors along
each auxiliary prefix of length $k$ and summing over prefixes proves
\eqref{eq:root-rn}.  Histories stopped before completing $k$ edges contribute
zero.  A previously visited non-root vertex cannot be the first repeat because
its target column was consumed when the path first entered it.  Before the
$j$th opportunity, exactly $j$ non-root targets have been removed from the
closure row.  When that row has width at most $h$,
Lemma~\ref{lem:predictable-exposure} gives \eqref{eq:root-hazard}.
\end{proof}

\begin{lemma}[Harmonic product]
\label{lem:harmonic-product}
Fix $\zeta\in(0,1)$.  Uniformly for $c\in[1/2,2]$, integers $m\ge1$, and
$0\le s\le(1-\zeta)m$,
\begin{equation}
 \prod_{j=0}^{s-1}\left(1-\frac c{m-j}\right)
 =(1-s/m)^c\{1+O_\zeta(m^{-1})\},
 \label{eq:harmonic-product}
\end{equation}
whenever all factors are nonnegative.  Also, uniformly for
$x\in[\zeta,1-\zeta]$ and $|a|\le1/2$,
\[
 |(1-x)^{1+a}-(1-x)|\le C_\zeta|a|.
\]
\end{lemma}

\begin{proof}
Expand the logarithm.  Since $m-j\ge\zeta m$,
\[
 \sum_{j<s}\log\left(1-\frac c{m-j}\right)
 =-c\sum_{j<s}\frac1{m-j}+O_\zeta(m^{-1}),
\]
and the harmonic sum equals
$\log(m/(m-s))+O_\zeta(m^{-1})$.  The second assertion follows from the
mean-value theorem in the exponent.
\end{proof}

\subsection{A uniform residual fraction}

\begin{proposition}[Uniform conditional residual fraction]
\label{prop:uniform-root-fraction}
Fix partition endpoints $0<S<U\le T$, $\xi>0$,
$\zeta\in(0,1/2)$, and $0\le\delta<1/4$.  On $\cG_n$, let $\cH$ be a
$(U,\delta,\xi)$-regular history.  Choose a root uniformly from its $N$
remaining vertices, independently of the unexposed matchings, and let $L$ be
the root-cycle length.  Uniformly over all such histories,
\begin{align}
 &\sup_{x\in[\zeta,1-\zeta]}
 \left|\Pp\left(\left.\frac LN\le x\right|\cH\right)-x\right|
 \notag\\
 &\qquad\le
 C_{S,\xi,\zeta}\Delta_n(U,\delta,h)
 +C_\xi p_n(D(S))+\pi_n+\frac{C_{\cJ,\xi,\zeta}}n.
 \label{eq:root-cdf}
\end{align}
\end{proposition}

\begin{proof}
Let $v$ be the uniform root and $C=c(v)$.  By
\eqref{eq:cycle-row-balance},
\[
 \Pp(C=a\mid\cH)=\frac{|Q_{n,a}\cap R|}{N}=\frac{m_a^R}{N}.
\]
Therefore
\begin{equation}
 \Pp(C\in D(S)\mid\cH)
 \le\frac{np_n(D(S))}{N}\le\xi^{-1}p_n(D(S)).
 \label{eq:closure-tail}
\end{equation}
Work henceforth on $C\notin D(S)$.  Put $m=m_C^R$ and $p_C=p_{n,C}$.
Regularity gives
\begin{equation}
 \frac{m}{Np_C}=1+O(\delta).
 \label{eq:closure-size}
\end{equation}
For the fixed partition, $p_C$ is bounded below over inner rows, so
$m\ge c_{\cJ,S}\xi n$.

First exclude premature termination of the avoid-$v$ exploration.  If it stops
after $r$ edges because only $v$ remains in the closure row, then
$R_C(r)=m-1$.  On the simultaneous event of
Lemma~\ref{lem:orbit-balance},
\[
 \left|\frac{m-1}{rp_C}-1\right|
 \le C_S\frac nr\Delta_n+\frac{C_{\cJ}}r.
\]
Multiplying by $r/N$ and using \eqref{eq:closure-size} gives
\begin{equation}
 |1-r/N|\le C_{S,\xi}\Delta_n+O_{\cJ,\xi}(n^{-1}).
 \label{eq:no-premature}
\end{equation}
Unless the desired bound is already trivial, this excludes termination before
$(1-\zeta)N$ edges.

Fix $x\in[\zeta,1-\zeta]$ and $k=\lfloor xN\rfloor$.  On the simultaneous
event above, $S(k)=R_C(k)$.  Lemma~\ref{lem:orbit-balance} and
\eqref{eq:closure-size} give
\begin{equation}
 \sup_{x\in[\zeta,1-\zeta]}
 \left|\frac{S(\lfloor xN\rfloor)}m-x\right|
 \le C_{S,\xi,\zeta}\Delta_n+O_{\cJ,\xi,\zeta}(n^{-1})
 \label{eq:closure-opportunities}
\end{equation}
except on conditional probability at most $\pi_n$.

On this event, the weight $W_k$ in Lemma~\ref{lem:root-avoiding} is squeezed between the
two harmonic products with $c=e^\beta$ and $c=e^{-\beta}$.  Since the mesh is
chosen with $\beta$ small, Lemma~\ref{lem:harmonic-product} and
\eqref{eq:closure-opportunities} show that both differ from $1-x$ by the
principal error in \eqref{eq:root-cdf}.  Average under $\Pp^\circ$, use the
trivial bound on the exceptional event, and add \eqref{eq:closure-tail}.
\end{proof}

\subsection{Deletion stability}

\begin{lemma}[Deletion of an interior macroscopic cycle]
\label{lem:deletion-stability}
Fix partition endpoints $0<V<U\le T$, $\xi>0$,
$\zeta\in(0,1/2)$, and $0\le\delta<1/4$.  On $\cG_n$, start from a
$(U,\delta,\xi)$-regular history and expose the cycle of a uniform residual
root.  Let $N$ and $L$ be the residual size and cycle length.  There is a
constant $C=C(V,\xi,\zeta)$ such that
\begin{align}
 &\Pp\Bigl(
 \zeta N\le L\le(1-\zeta)N,\ 
 \text{the post-deletion history is not }
 (V,\delta',\zeta\xi)\text{-regular}
 \mid\cH\Bigr)
 \le\pi_n,
 \label{eq:deletion-joint}
\end{align}
where
\begin{equation}
 \delta'=C\Delta_n(U,\delta,h)+\frac{C_{\cJ,\xi,\zeta}}n.
 \label{eq:deletion-error}
\end{equation}
The assertion is used only when $\delta'<1/4$.
\end{lemma}

\begin{proof}
Let $d_a$ be the number of vertices of the exposed cycle in $H_{n,a}$ and
$d_{ab}$ the number of its edges from $H_{n,a}$ to $H_{n,b}$.  The event in
Lemma~\ref{lem:orbit-balance} is simultaneous over all prefix lengths and may
therefore be evaluated at the random closing time without conditioning on
$L$.  On $\zeta N\le L\le(1-\zeta)N$, apply
\eqref{eq:scale-sensitive-balance} with inner cutoff $V$ and $\ell=L$ to get
\[
 \frac{d_a}{Lp_{n,a}}
 =1+O_{V,\xi,\zeta}(\Delta_n)+O_{\cJ,\xi,\zeta}(n^{-1})
\]
for every row below $V$.  Since
$m_a^R=Np_{n,a}(1+O(\delta))$,
\begin{equation}
 \frac{d_a}{m_a^R}
 =\frac LN+O_{V,\xi,\zeta}(\Delta_n)+O_{\cJ,\xi,\zeta}(n^{-1}).
 \label{eq:deleted-row-fraction}
\end{equation}
The row-prefix estimate at the same closing time gives
\[
 d_{ab}=\frac{d_a}{m_a^R}N_{ab}^R
        +O((\beta+\chi_n)N_{ab}^R),
\]
so
\begin{equation}
 d_{ab}=\frac LN N_{ab}^R
 +\bigl\{O_{V,\xi,\zeta}(\Delta_n)
          +O_{\cJ,\xi,\zeta}(n^{-1})\bigr\}N_{ab}^R.
 \label{eq:deleted-cell-fraction}
\end{equation}
The quantity $d_{ab}$ is exactly the number of targets removed from
$Q_{n,a}\cap H_{n,b}$.  Equations
\eqref{eq:deleted-row-fraction}--\eqref{eq:deleted-cell-fraction}, together
with $1-L/N\ge\zeta$, prove regularity of the residual history with error
\eqref{eq:deletion-error}.  Residual factorization remains valid by
Lemma~\ref{lem:residual-factorization}.
\end{proof}

\section{The regeneration induction}
\label{sec:regeneration}

Define residual fractions
\begin{equation}
 V_{n,j}=\frac{\widetilde L_{n,j}}
 {n-\widetilde L_{n,1}-\cdots-\widetilde L_{n,j-1}},
 \label{eq:residual-fractions}
\end{equation}
with $V_{n,j}=0$ after exhaustion.

\begin{lemma}[Conditional CDF test]
\label{lem:cdf-test}
Let $W\in[0,1]$, let $\mathscr F$ be a sigma-field, and suppose
\[
 \sup_{x\in[a,b]}|\Pp(W\le x\mid\mathscr F)-x|\le\eta
 \quad\text{a.s.},
 \qquad 0<a<b<1.
\]
Then
\[
 \Pp(W\notin[a,b]\mid\mathscr F)\le a+(1-b)+2\eta.
\]
If $f$ is absolutely continuous, supported in $[a,b]$, and vanishes at the
endpoints, then
\[
 \left|\Ee(f(W)\mid\mathscr F)-\int_0^1f(u)\dd u\right|
 \le\eta\int_a^b|f'(u)|\dd u.
\]
\end{lemma}

\begin{proof}
The endpoint estimate follows by evaluating the CDF bound at $a,b$.  The second
claim is conditional Stieltjes integration by parts.
\end{proof}

\begin{proposition}[Independent residual uniforms]
\label{prop:residual-uniforms}
For every fixed $k$,
\begin{equation}
  (V_{n,1},\ldots,V_{n,k})
  \Rightarrow(U_1,\ldots,U_k),
  \qquad U_j\stackrel{\mathrm{iid}}\sim\Unif(0,1).
  \label{eq:finite-uniforms}
\end{equation}
\end{proposition}

\begin{proof}
Fix $k$, an interior parameter $\zeta\in(0,1/4)$, and a target error
$\eta\in(0,1/16)$.  Put $\xi_*=\zeta^{k-1}$.  We describe the order of
choices, which is essential.

Choose an inner cutoff $T_0$ so large that
$\xi_*^{-1}\overline\nu(T_0)\le\eta/4$.  Once $T_j$ is fixed, choose
constants $R_j,D_j\ge1$ bounding the principal-error coefficients in,
respectively, Proposition~\ref{prop:uniform-root-fraction} and
Lemma~\ref{lem:deletion-stability}, with inner cutoff $T_j$, residual fraction
at least $\xi_*$, and truncation parameter $\zeta$.  These constants are
independent of the outer cutoff and of the subsequent partition.  Set
\[
 A_j=\prod_{i=0}^{j}(1+R_i+D_i),
\]
and choose $T_{j+1}>T_j$ so large that
\[
 A_j\sqrt{\overline\nu(T_{j+1})}\le2^{-j-3}\eta.
\]
Continue through $T_k$.  This is possible because
$\overline\nu(U)\downarrow0$.

Write $\epsilon_{j+1}=\sqrt{\overline\nu(T_{j+1})}$.  An error introduced at
level $j$ contributes to a later regularity error at level $i\le j$ with
coefficient $\prod_{\ell=i}^jD_\ell$, and to the root-law error at level $i$
with coefficient $R_i\prod_{\ell=i+1}^jD_\ell$, with empty products equal to
one.  Both coefficients are at most $A_j$.  Thus the total tail contribution
to any regularity or root-law error is at most
\[
 \sum_{j=0}^{k-1}A_j\epsilon_{j+1}\le\eta/4.
\]

After all cutoffs and constants are fixed, choose $h>0$ such that
\[
 2h\le\beta_0,
 \qquad
 2h\sum_{j=0}^{k-1}A_j\le\eta/4.
\]
Choose one finite partition of $[0,T_k]$ containing $T_0,\ldots,T_k$ and
having mesh at most $h$, and append its tail cell.  The same propagation
calculation bounds every mesh contribution by $\eta/4$.  Finally take $n$
large enough that the fixed-partition errors, the empirical tail-mass errors,
and the exceptional-event probabilities, after these finitely many
amplifications, contribute at most another prescribed fraction of $\eta$.
Every required regularity error is then below $1/4$, and every stagewise
conditional CDF error is $O_k(\eta)$ without a cutoff-dependent factor
multiplying $\eta$.  For reference, at deletion stage $1\le r\le k$ we use
\[
\begin{array}{c|c|c|c}
\text{stage}&\text{current outer cutoff}&\text{root-law inner cutoff}
&\text{post-deletion cutoff}\\ \hline
r&T_{k-r+1}&T_{k-r}&T_{k-r}.
\end{array}
\]
Thus every inner cutoff is fixed before its corresponding outer cutoff.

On the event
\begin{equation}
 \zeta\le V_{n,j}\le1-\zeta,
 \qquad 1\le j<r,
 \label{eq:interior-history}
\end{equation}
the residual size before stage $r$ is at least
$\zeta^{r-1}n\ge\xi_*n$.  Inductively, Lemma~\ref{lem:deletion-stability}
keeps the current history regular at the required outer cutoff, apart from a
conditional probability $O_k(\eta)$.  Proposition
\ref{prop:uniform-root-fraction} then gives
\begin{equation}
 \sup_{x\in[\zeta,1-\zeta]}
 \left|\Pp(V_{n,r}\le x\mid\mathscr F_{r-1})-x\right|
 \le C_k\eta
 \label{eq:stage-cdf}
\end{equation}
uniformly over every surviving regular history.

Let $f_1,\ldots,f_k$ be bounded Lipschitz functions on $[0,1]$.  Choose a
piecewise linear cutoff $\varphi_\zeta$ supported in
$[\zeta,1-\zeta]$ and equal to one on
$[2\zeta,1-2\zeta]$, and put $g_j=f_j\varphi_\zeta$.  Apply
Lemma~\ref{lem:cdf-test} to \eqref{eq:stage-cdf}, multiply successively by the
preceding factors, and use the joint-event form of
Lemma~\ref{lem:deletion-stability}.  Induction gives
\[
 \Ee\prod_{j=1}^kg_j(V_{n,j})
 =\prod_{j=1}^k\int_0^1g_j(u)\dd u+O_{f,k,\zeta}(\eta)+o_n(1).
\]
The endpoint part of Lemma~\ref{lem:cdf-test} gives
\[
 \limsup_{n\to\infty}\Pp\left(
 V_{n,j}\notin[2\zeta,1-2\zeta]
 \text{ for some }j\le k\right)
 \le4k\zeta+C_k\eta.
\]
Replace $g_j$ by $f_j$, then let $n\to\infty$, $\eta\downarrow0$, and
$\zeta\downarrow0$.  Products of bounded Lipschitz functions determine weak
convergence on $[0,1]^k$, proving \eqref{eq:finite-uniforms}.
\end{proof}

\begin{proof}[Proof of Theorem~\ref{thm:main-spectral}]
Put
\begin{equation}
 X_{n,j}=\frac{\widetilde L_{n,j}}n
 =V_{n,j}\prod_{i<j}(1-V_{n,i}),
 \qquad
 R_{n,k}=1-\sum_{j=1}^kX_{n,j}
 =\prod_{i=1}^k(1-V_{n,i}).
 \label{eq:stick-breaking-finite}
\end{equation}
Proposition~\ref{prop:residual-uniforms} and the continuous mapping theorem give
for every fixed $k$ the corresponding finite-dimensional $\GEM(1)$ limit,
including
\[
 R_{n,k}\Rightarrow\prod_{i=1}^k(1-U_i).
\]
Since $0\le R_{n,k}\le1$,
\begin{equation}
 \Ee R_{n,k}\longrightarrow2^{-k}.
 \label{eq:residual-expectation}
\end{equation}
Let $X_n=(X_{n,1},X_{n,2},\ldots)$ and let $X$ have the $\GEM(1)$ law.  If
$F:\ell^1\to\mathbb R$ is bounded and Lipschitz, truncation after $k$
coordinates gives
\[
 |\Ee F(X_n)-\Ee F(X)|
 \le\operatorname{Lip}(F)(\Ee R_{n,k}+2^{-k})
 +|\Ee F(X_n^{(k)})-\Ee F(X^{(k)})|.
\]
First let $n\to\infty$ and then $k\to\infty$.  The bounded-Lipschitz criterion
on the separable space $\ell^1$ proves \eqref{eq:gem-main}.

For nonnegative $\ell^1$ sequences, decreasing rearrangement is a contraction:
\[
 \norm{x^\downarrow-y^\downarrow}_1\le\norm{x-y}_1.
\]
This follows first for finite sequences by the two-point uncrossing inequality
and then by truncation.  The continuous mapping theorem proves
\eqref{eq:pd-main}.  The first residual fraction is the cycle length of a
uniform root divided by $n$, proving \eqref{eq:uniform-root-main}.

The argument was carried out along an arbitrary subsequence after extracting a
marked limit.  Lemma~\ref{lem:channel-domination} gives the same strict bound
$q(\nu)$ for every such marked limit, and the resulting law is always universal.
The subsequence principle therefore yields convergence along the full sequence.
\end{proof}

\begin{corollary}[Power sums]
\label{cor:power-sums}
Under the hypotheses of Theorem~\ref{thm:main-spectral}, for every $s>1$,
\[
 \sum_{j\ge1}\left(\frac{L_{n,j}}n\right)^s
 \Rightarrow\sum_{j\ge1}P_j^s,
 \qquad (P_j)\sim\PD(1),
\]
and
\[
 \Ee\sum_{j\ge1}\left(\frac{L_{n,j}}n\right)^s\longrightarrow\frac1s.
\]
\end{corollary}

\begin{proof}
The power-sum functional is bounded and continuous on $\cS^\downarrow$.  A
size-biased pick $\widetilde P$ from a $\PD(1)$ partition is uniform on
$(0,1)$, and
\[
 \Ee(\widetilde P^{s-1}\mid(P_j))=\sum_jP_j^s.
\]
Hence the limiting expectation is
$\int_0^1u^{s-1}\dd u=1/s$.
\end{proof}

\section{Regularly varying spectral gaps}
\label{sec:rv-gap}

We prove Theorem~\ref{thm:rv-gap}.  Let
\[
  F(x)=\nu((0,x])
\]
be regularly varying at zero with index $\alpha>0$.  Define the truncated first
moment
\begin{equation}
  H(x)=\int_{(0,x]}r\nu(\dd r).
  \label{eq:H-first-moment}
\end{equation}
Integration by parts and Karamata's theorem give
\begin{equation}
  H(x)=xF(x)-\int_0^xF(u)\dd u
  \sim\frac{\alpha}{\alpha+1}xF(x)
  \qquad(x\downarrow0).
  \label{eq:H-asymptotic}
\end{equation}
Thus $H$ is regularly varying at zero with index $\alpha+1$.  Since
$B_\nu(t)=\int e^{-tr}\dd H(r)$, the Abelian theorem for Laplace--Stieltjes
transforms yields
\begin{equation}
  B_\nu(t)
  \sim \Gamma(\alpha+2)H(1/t)
  \sim \alpha\Gamma(\alpha+1)t^{-1}F(1/t)
  \qquad(t\to\infty).
  \label{eq:B-rv-asymptotic}
\end{equation}
These standard forms of Karamata's theorem are recorded in
\cite[Chapters~1 and~4]{BGT1987}.

Let $C=C_\nu$ and $Q=CC^*$ on $L^2(\nu)$.  A direct calculation gives the
symmetric Markov kernel
\begin{equation}
  R_\nu(r,s)=rs\int_0^\infty
  \frac{e^{-(r+s)t}}{B_\nu(t)}\dd t
  \label{eq:Rnu}
\end{equation}
relative to $\nu(\dd s)$.

\begin{lemma}[Dyadic kernel lower bound]
\label{lem:rv-kernel-lower}
Let
\[
 I_k=(2^{-k-1},2^{-k}],
 \qquad w_k=\nu(I_k).
\]
There are $K<\infty$ and $c>0$ such that, for every $k\ge K$, if $r,s$ lie in
the same shell $I_k$ or in the adjacent shells $I_k,I_{k-1}$, then
\begin{equation}
  R_\nu(r,s)\ge\frac{c}{w_k}.
  \label{eq:rv-shell-lower}
\end{equation}
Moreover,
\begin{equation}
  \frac{w_{k+1}}{w_k}\longrightarrow2^{-\alpha}.
  \label{eq:shell-ratio}
\end{equation}
\end{lemma}

\begin{proof}
Put $x=2^{-k}$.  Regular variation gives
\[
  w_k=F(x)-F(x/2)\sim(1-2^{-\alpha})F(x),
\]
which proves \eqref{eq:shell-ratio}.  By
\eqref{eq:B-rv-asymptotic} and the uniform convergence theorem, uniformly for
$t\in[x^{-1},2x^{-1}]$ and all sufficiently small $x$,
\begin{equation}
  B_\nu(t)\le CxF(x).
  \label{eq:B-shell-upper}
\end{equation}
If $r,s$ are in the stated shells, then $rs\ge c_1x^2$ and
$r+s\le c_2x$.  Restricting \eqref{eq:Rnu} to
$t\in[x^{-1},2x^{-1}]$ therefore gives
\[
 R_\nu(r,s)
 \ge c x^2\int_{x^{-1}}^{2x^{-1}}
       \frac{\dd t}{xF(x)}
 \ge\frac{c'}{F(x)}
 \ge\frac{c''}{w_k}.
\]
\end{proof}

We also need a weighted Poincar\'e inequality for an eventually geometric
sequence of shell masses.

\begin{lemma}[Geometric weighted Poincar\'e inequality]
\label{lem:geometric-poincare}
Let $(w_k)_{k\ge K}$ be positive and suppose that
$w_{k+1}/w_k\le\rho<1$ for all $k\ge K$.  For a sequence $(m_k)$,
\begin{equation}
 \sum_{k\ge K}w_k(m_k-m_K)^2
 \le C_\rho\sum_{k>K}w_k(m_k-m_{k-1})^2.
 \label{eq:geometric-poincare}
\end{equation}
Consequently the same right-hand side controls the weighted variance of
$(m_k)_{k\ge K}$.
\end{lemma}

\begin{proof}
Put $d_j=m_j-m_{j-1}$ and $a_j=\sqrt{w_j}d_j$.  For $j\le k$,
$w_k/w_j\le\rho^{k-j}$, and hence
\[
 \sqrt{w_k}(m_k-m_K)
 =\sum_{j=K+1}^k\sqrt{\frac{w_k}{w_j}}a_j.
\]
Young's convolution inequality on $\ell^2$, with the summable kernel
$(\rho^{r/2})_{r\ge0}$, proves \eqref{eq:geometric-poincare}.  Weighted
variance is no larger than squared distance from the fixed value $m_K$.
\end{proof}

\begin{proof}[Proof of Theorem~\ref{thm:rv-gap}]
Choose $K$ so that Lemma~\ref{lem:rv-kernel-lower} holds and
$w_{k+1}/w_k\le\rho<1$ for $k\ge K$.  Put
$I_\circ=(2^{-K},1]$ and $w_\circ=\nu(I_\circ)$.  For $K$ large,
$w_\circ>0$ and $w_K>0$.  Restricting \eqref{eq:Rnu} to $t\in[0,1]$, and
using $B_\nu(t)\le1$, gives a positive lower bound for $R_\nu$ on
$[2^{-K-1},1]^2$.  In particular, the Dirichlet form controls both the variance
inside the core and the squared difference between the core mean and the mean
on $I_K$.

Because $Q$ is a symmetric Markov operator with kernel $R_\nu$ relative to
$\nu$, its Dirichlet form is
\begin{equation}
 \langle f,(I-Q)f\rangle
 =\frac12\iint (f(r)-f(s))^2R_\nu(r,s)\,\nu(\dd r)\nu(\dd s).
 \label{eq:rv-dirichlet-form}
\end{equation}
Let $m_\circ$ be the $\nu$-average of $f$ on $I_\circ$ and let $m_k$ be its
average on $I_k$.  The same-shell part of
\eqref{eq:rv-shell-lower} yields
\begin{equation}
 \langle f,(I-Q)f\rangle
 \ge c\left\{\int_{I_\circ}|f-m_\circ|^2\dd\nu
       +\sum_{k\ge K}\int_{I_k}|f-m_k|^2\dd\nu\right\},
 \label{eq:rv-within-shell}
\end{equation}
while the core--first-shell and adjacent-shell parts, followed by Jensen's
inequality, give
\begin{equation}
 \langle f,(I-Q)f\rangle
 \ge c\left\{w_K(m_K-m_\circ)^2
       +\sum_{k>K}w_k(m_k-m_{k-1})^2\right\}.
 \label{eq:rv-between-shell}
\end{equation}
Since $\sum_{k\ge K}w_k\le w_K/(1-\rho)$,
Lemma~\ref{lem:geometric-poincare} gives
\begin{align}
 \sum_{k\ge K}w_k(m_k-m_\circ)^2
 &\le 2C_\rho\sum_{k>K}w_k(m_k-m_{k-1})^2
       +\frac{2w_K}{1-\rho}(m_K-m_\circ)^2.
 \label{eq:core-shell-poincare}
\end{align}
Combining this with \eqref{eq:rv-within-shell}--\eqref{eq:rv-between-shell}
and $\Var_\nu(f)\le\int|f-m_\circ|^2\dd\nu$ yields
\begin{equation}
  \Var_\nu(f)
  \le C_\nu\langle f,(I-Q)f\rangle.
  \label{eq:rv-poincare}
\end{equation}
Since $Q=CC^*$ is positive semidefinite, \eqref{eq:rv-poincare} implies
$\norm{Q}_{L^2_0(\nu)}\le1-c_\nu$ for some $c_\nu>0$.  Therefore
\[
  q(\nu)=\norm{C}_{L^2_0(\mu_\nu)\to L^2_0(\nu)}
  \le\sqrt{1-c_\nu}<1.
\]
The power-law law \eqref{eq:power-law-law} has $F(x)=x^\alpha$, so it is a
special case.
\end{proof}

\begin{remark}
The proof gives no sharp value of $q(\nu)$.  Even for the exact power-law law,
determining the optimal gap appears to be a nontrivial Mellin-analytic problem.
\end{remark}

\section{Applications}
\label{sec:applications}

\subsection{Comparable arrays}

\begin{corollary}[Comparable rates]
\label{cor:comparable}
Suppose that after common rescaling all rates lie in $[c,1]$ for some fixed
$c>0$.  Then the conclusions of Theorem~\ref{thm:main-spectral} hold.
\end{corollary}

\begin{proof}
For every $r\in[c,1]$,
\[
 r e^{-rt}\dd t\ge c e^{-t}\dd t.
\]
Thus all conditional time laws contain the common probability component
$e^{-t}\dd t$ with mass $c$.  The common-component decomposition and
Jensen's inequality give $q(\nu)\le\sqrt{1-c}<1$ for every subsequential
empirical rate law $\nu$.  Although the full sequence of empirical laws need
not converge, every subsequence has a further convergent subsequence, and
Theorem~\ref{thm:main-spectral} gives the same universal limit on each such
refinement.  The subsequence principle proves the full conclusion.
\end{proof}

\subsection{Power-law and Sukhatme spectra}

\begin{proof}[Proof of Corollary~\ref{cor:power-law-spectra}]
Normalize by $n^\gamma$.  The empirical measure of
$(j/n)^\gamma$, $1\le j\le n$, converges to the pushforward of Lebesgue measure
under $x\mapsto x^\gamma$.  Its density is
\[
 \frac1\gamma r^{1/\gamma-1}\dd r=\nu_{1/\gamma}(\dd r).
\]
The assignment of this multiset to labels does not change the empirical rate
law.  Apply Theorems~\ref{thm:main-spectral} and
\ref{thm:rv-gap}.  The Sukhatme case is $\gamma=1$.
\end{proof}

The ordered Sukhatme profile $n,n-1,\ldots,1$ also has a distinct microscopic
terminal theory: its fixed points have logarithmic intensity and a terminal
Poisson process; see \cite{LongSukhatme2026} and the historical exponential
spacing representations of Sukhatme~\cite{Sukhatme1937} and
R\'enyi~\cite{Renyi1953}.  The present macroscopic result is insensitive to
ordering and therefore belongs to a different universality mechanism.

\subsection{Regularly varying spectra}

\begin{lemma}[Empirical law of a regularly varying spectrum]
\label{lem:rv-empirical}
Let $a_m=m^\gamma L(m)$ be positive and regularly varying with index
$\gamma>0$, and put $M_n=\max_{1\le j\le n}a_j$.  Then
\[
 \frac{M_n}{a_n}\longrightarrow1
\]
and
\begin{equation}
 \frac1n\sum_{j=1}^n\delta_{a_j/M_n}
 \Rightarrow\nu_{1/\gamma}.
 \label{eq:rv-empirical}
\end{equation}
\end{lemma}

\begin{proof}
For every fixed $\delta\in(0,1)$, the uniform convergence theorem for regular
variation gives
\[
 \sup_{x\in[\delta,1]}
 \left|\frac{a_{\lfloor nx\rfloor}}{a_n}-x^\gamma\right|\to0.
\]
Potter bounds imply that for any $\epsilon\in(0,\gamma)$ there is $C_\epsilon$
such that, apart from finitely many indices,
\[
 \frac{a_j}{a_n}\le C_\epsilon(j/n)^{\gamma-\epsilon},
 \qquad j\le n.
\]
Since $M_n\ge a_n$, it remains to prove the matching upper bound.  On
$[\delta n,n]$ this follows from uniform convergence, while Potter bounds make
$\sup_{j\le\delta n}a_j/a_n$ arbitrarily small after first choosing $\delta$
and then $n$ large; the finitely many exceptional initial indices are
$o(a_n)$.  Hence $M_n/a_n\to1$.  For every bounded continuous $\varphi$, split
the Riemann sum at $\delta n$.  The small-index part has mass at most $\delta$, while on
$[\delta,1]$ uniform convergence gives
\[
 \frac1n\sum_{j=1}^n\varphi(a_j/M_n)
 \longrightarrow\int_0^1\varphi(x^\gamma)\dd x.
\]
The latter is integration against $\nu_{1/\gamma}$.
\end{proof}

\begin{proof}[Proof of Corollary~\ref{cor:regularly-varying}]
The empirical normalized rate law is independent of the assignment to labels.
Lemma~\ref{lem:rv-empirical} identifies its limit as
$\nu_{1/\gamma}$, and Theorems~\ref{thm:main-spectral} and
\ref{thm:rv-gap} apply.
\end{proof}

For the exact power-law spectrum, the late-clock mass has the explicit decay
\begin{equation}
 \overline\nu(U)
 =\alpha\int_0^1r^{\alpha-1}e^{-rU}\dd r
 \asymp U^{-\alpha}.
 \label{eq:power-tail}
\end{equation}
Thus the uncontrolled-row term in orbit balance is
$O(U^{-\alpha/2})=O(U^{-1/(2\gamma)})$.  This is qualitative here because the
mesh and cutoffs are fixed before $n\to\infty$, but it identifies the correct
endpoint leakage scale.

\section{Spectral boundary and further questions}
\label{sec:boundary}

The spectral formulation separates three inputs: vanishing late-clock mass,
local flatness within finite clock slabs, and a global mean-zero contraction.
We first establish a broader sufficient condition for the channel gap and
identify its slowly varying boundary.  The necessity of a gap for
Poisson--Dirichlet universality and the behavior of multiple rate scales remain
separate questions.

\begin{theorem}[Dyadic-doubling criterion]
\label{thm:dyadic-doubling}
Let $\nu$ be a probability law on $(0,1]$, and put
\[
 w_k=\nu((2^{-k-1},2^{-k}]).
\]
Suppose that, for some constants $0<a\le b<1$ and every sufficiently large
$k$, one has $w_k>0$ and
\[
 a\le\frac{w_{k+1}}{w_k}\le b.
\]
Then $q(\nu)<1$.
\end{theorem}

\begin{proof}
Choose $K$ so that the ratio bounds hold for $k\ge K$, and let
$I_k=(2^{-k-1},2^{-k}]$.  Put $x=2^{-k}$ and
$t\in[x^{-1},2x^{-1}]$.  The contribution of the inner shells to
$B_\nu(t)$ is at most
\[
 \sum_{m\ge0}2^{-k-m}w_{k+m}
 \le xw_k\sum_{m\ge0}(b/2)^m.
\]
For an outer shell $I_{k-m}$ with $1\le m\le k-K$, one has
$r\le2^mx$, $rt\ge2^{m-1}$, and $w_{k-m}\le a^{-m}w_k$.  Its
contribution is at most
\[
 xw_k(2/a)^m e^{-2^{m-1}}.
\]
The sum over $m\ge1$ converges.  The remaining core $(2^{-K},1]$ contributes
at most $e^{-2^{-K}/x}=o(xw_k)$, since $w_k\ge w_Ka^{k-K}$.  Thus, uniformly
on the stated time interval and for all sufficiently large $k$,
\begin{equation}
 B_\nu(t)\le Cxw_k.
 \label{eq:dyadic-density-upper}
\end{equation}

If $r,s$ lie in $I_k$, or in the adjacent shells $I_k,I_{k-1}$, then
$rs\ge c_1x^2$ and $r+s\le c_2x$.  Restricting \eqref{eq:Rnu} to
$[x^{-1},2x^{-1}]$ and using \eqref{eq:dyadic-density-upper} gives
\[
 R_\nu(r,s)\ge c/w_k.
\]
After enlarging $K$ if necessary, the core has positive mass.  Restricting
\eqref{eq:Rnu} to $t\in[0,1]$ gives a positive kernel lower bound joining
the core to $I_K$ and within the core.  The within-shell and adjacent-shell
estimates in \eqref{eq:rv-within-shell}--\eqref{eq:rv-between-shell} therefore
hold in this setting.  Lemma~\ref{lem:geometric-poincare} and
\eqref{eq:core-shell-poincare}, with $\rho=b$, give
\[
 \Var_\nu(f)\le C_\nu\langle f,(I-C_\nu C_\nu^*)f\rangle.
\]
Positivity of $C_\nu C_\nu^*$ now yields $q(\nu)<1$.
\end{proof}

The criterion permits oscillating shell ratios and atomic or singular rate
laws.  It also implies Theorem~\ref{thm:rv-gap}, because a regularly varying
lower tail of index $\alpha>0$ has $w_{k+1}/w_k\to2^{-\alpha}\in(0,1)$.

\begin{theorem}[Slowly varying boundary]
\label{thm:slow-boundary}
Let $\nu$ be a probability law on $(0,1]$.  Suppose
$F_\nu(x)=\nu((0,x])$ is positive for all sufficiently small $x$, tends to
zero, and is slowly varying at zero.
Then $q(\nu)=1$.
\end{theorem}

\begin{proof}
In the joint exponential channel, define
\[
 A_\delta=\{R\le\delta\},\qquad
 B_\delta=\{T>1/\delta\},\qquad
 a_\delta=F_\nu(\delta).
\]
The Laplace--Stieltjes Abelian theorem at index zero
\cite[Chapter~1]{BGT1987} gives
\[
 \Pp(B_\delta)=\int e^{-r/\delta}\nu(\dd r)\sim a_\delta.
\]
Also, $1-e^{-z}\le z$ and Stieltjes integration by parts yield
\begin{align*}
 \Pp(A_\delta\setminus B_\delta)
 &\le\delta^{-1}\int_{(0,\delta]}r\nu(\dd r)\\
 &=F_\nu(\delta)-\delta^{-1}\int_0^\delta F_\nu(u)\dd u
 =o(a_\delta).
\end{align*}
For the last equality, substitute $u=\delta v$ and apply dominated
convergence to $F_\nu(\delta v)/F_\nu(\delta)$ on $0<v<1$: slow variation
gives pointwise convergence to one, and monotonicity bounds the ratio by one.
Consequently $\Pp(A_\delta\cap B_\delta)\sim a_\delta$, while both marginal
probabilities are asymptotic to $a_\delta\to0$.  Hence
\[
 \Corr(\one_{A_\delta},\one_{B_\delta})\longrightarrow1.
\]
The channel norm $q(\nu)$ dominates each of these correlations and is at most
one, proving the claim.
\end{proof}

Theorem~\ref{thm:slow-boundary} identifies a boundary of the strict spectral-gap
hypothesis.  It does not assert failure of the $\PD(1)$ conclusion in that
regime; necessity of the gap remains the following question.

\begin{question}[Necessity of the spectral gap]
\label{q:sharpness}
Is $q(\nu)<1$ necessary for $\PD(1)$ universality uniformly over all
deterministic arrangements with empirical rate law $\nu$?  If $q(\nu)=1$,
can one construct an arrangement whose macroscopic cycle partition is not
$\PD(1)$?
\end{question}

The present proof genuinely loses regeneration when the mean-zero resolvent is
unbounded.  It is unknown whether that reflects a transition of the model or
only a limitation of the method.

A useful continuum-of-scales test family is explicit.  For $\beta>0$, let
\begin{equation}
 \nu_\beta(\dd r)=\frac{1}{\beta r}
 \one_{[e^{-\beta},1]}(r)\dd r.
 \label{eq:log-uniform-law}
\end{equation}
This is the limiting normalized rate law of
$\theta_{n,i}=\exp(\beta i/n)$.

\begin{proposition}[Degeneration of the log-uniform channel]
\label{prop:log-uniform-lower}
For every $\beta>0$,
\begin{equation}
 q(\nu_\beta)\ge
 \frac{\beta}{\sqrt{\beta^2+2\pi^2}}
 =1-\frac{\pi^2}{\beta^2}+O(\beta^{-4}).
 \label{eq:log-uniform-gap-lower}
\end{equation}
\end{proposition}

\begin{proof}
Write $R=e^{-X}$ with $X\sim\Unif(0,\beta)$.  If $E\sim\Exp(1)$ is
independent, then $T=E/R=Ee^X$, so
$\log T=X+\log E$.  Maximal correlation dominates the ordinary correlation of
the centered functions $X=X(R)$ and $\log T$.  Since
$\Var X=\beta^2/12$ and $\Var(\log E)=\pi^2/6$,
\[
 q(\nu_\beta)\ge\Corr(X,\log T)
 =\sqrt{\frac{\Var X}{\Var X+\Var(\log E)}}
 =\frac{\beta}{\sqrt{\beta^2+2\pi^2}}.
\]
\end{proof}

Since $\nu_\beta$ is supported on $[e^{-\beta},1]$, the comparable-support
bound gives $q(\nu_\beta)<1$.  Thus, for every $\beta>0$, the norm gap satisfies
\[
 0<1-q(\nu_\beta)
 \le1-\frac{\beta}{\sqrt{\beta^2+2\pi^2}}
 \le\frac{\pi^2}{\beta^2}.
\]
The last inequality follows from $(1+x)^{-1/2}\ge1-x/2$ for $x\ge0$.

\begin{question}[Growing log-linear spectra]
\label{q:log-linear-growth}
For
\[
 \theta_{n,i}=\exp(\beta_n i/n),
 \qquad \beta_n\longrightarrow\infty,
\]
determine the largest-cycle scale and the maximal growth regime of $\beta_n$
for which the $\PD(1)$ limit persists.  Is there a sharp transition between
macroscopic and submacroscopic cycles?
\end{question}

After normalization the empirical spectrum collapses to an atom at zero, so
this problem is outside Theorem~\ref{thm:main-spectral}.  Unlike a finite
hierarchy of separated scales, it has a continuum of logarithmic rate scales.
It is therefore a natural test of whether one-scale spectral universality can
survive a slowly degenerating gap.

\begin{conjecture}[Multiscale endpoint phase]
\label{conj:multiscale}
Suppose that, after normalization by the largest rate scale, the empirical rate
law has a positive atom at zero, and that the vanishing-rate sector admits one
or more nontrivial limits after successive rescalings.  Then the cycle
partition has a hierarchical multiscale limit assembled from the
Poisson--Dirichlet partitions of the successive rate scales, with a
continuum-scale analogue when the logarithmic spectrum does not separate into
finitely many levels.
\end{conjecture}

An atom at zero is excluded from Theorem~\ref{thm:main-spectral} because a
positive fraction of clocks escapes every fixed time cutoff.  The conjecture
asks for the correct replacement rather than a direct extension of the
one-scale theorem.

\begin{question}[Quantitative convergence]
\label{q:quantitative}
Under a uniform spectral gap and a polynomial bound
$\overline\nu(U)\le CU^{-\alpha}$, can one obtain an explicit rate of
convergence to $\GEM(1)$ or $\PD(1)$ in Wasserstein or bounded-Lipschitz
metric?
\end{question}

Such a theorem would require growing partitions and quantitative approximation
of the finite-state spectral gap.  The present proof deliberately fixes all
cutoffs and the mesh before sending $n\to\infty$.

\begin{question}[Exact power-law gap]
\label{q:exact-gap}
Determine the exact value, or sharp asymptotics as $\alpha\downarrow0$ and
$\alpha\to\infty$, of $q(\nu_\alpha)$.  Is there a Mellin transform that
diagonalizes an appropriate scale-invariant comparison channel?
\end{question}

\begin{thebibliography}{99}

\bibitem{BCD2026}
J. Borga, S. Chatterjee, and P. Diaconis.
\newblock Permuton and local limits for the Luce model.
\newblock \emph{Annals of Applied Probability}, to appear.
\newblock \href{https://arxiv.org/abs/2509.07729v3}{arXiv:2509.07729v3}.

\bibitem{BGT1987}
N. H. Bingham, C. M. Goldie, and J. L. Teugels.
\newblock \emph{Regular Variation}.
\newblock Encyclopedia of Mathematics and its Applications, vol. 27.
\newblock Cambridge University Press, Cambridge, 1987.

\bibitem{Billingsley1999}
P. Billingsley.
\newblock \emph{Convergence of Probability Measures}, second edition.
\newblock Wiley, New York, 1999.

\bibitem{CDK2024}
S. Chatterjee, P. Diaconis, and G. B. Kim.
\newblock Enumerative theory for the Tsetlin library.
\newblock \emph{Journal of Algebra} 655 (2024), 139--162.
\newblock \href{https://doi.org/10.1016/j.jalgebra.2023.08.009}
{doi:10.1016/j.jalgebra.2023.08.009}.

\bibitem{Dudley2002}
R. M. Dudley.
\newblock \emph{Real Analysis and Probability}, second edition.
\newblock Cambridge University Press, Cambridge, 2002.

\bibitem{Durrett2019}
R. Durrett.
\newblock \emph{Probability: Theory and Examples}, fifth edition.
\newblock Cambridge University Press, Cambridge, 2019.

\bibitem{Ewens1972}
W. J. Ewens.
\newblock The sampling theory of selectively neutral alleles.
\newblock \emph{Theoretical Population Biology} 3 (1972), no.~1, 87--112.
\newblock \href{https://doi.org/10.1016/0040-5809(72)90035-4}
{doi:10.1016/0040-5809(72)90035-4}.

\bibitem{Guerder2026}
A. Guerder.
\newblock Cycle structure of random standardized permutations.
\newblock Preprint, 2026.
\newblock \href{https://arxiv.org/abs/2603.24127}{arXiv:2603.24127}.

\bibitem{Kallenberg2002}
O. Kallenberg.
\newblock \emph{Foundations of Modern Probability}, second edition.
\newblock Springer-Verlag, New York, 2002.

\bibitem{Kingman1975}
J. F. C. Kingman.
\newblock Random discrete distributions.
\newblock \emph{Journal of the Royal Statistical Society, Series B}
37 (1975), 1--22.
\newblock \href{https://doi.org/10.1111/j.2517-6161.1975.tb01024.x}
{doi:10.1111/j.2517-6161.1975.tb01024.x}.

\bibitem{LongPD2026}
C. D. Long.
\newblock Poisson--Dirichlet universality for Luce permutations.
\newblock Preprint, 2026.
\newblock Repository manuscript:
\href{https://github.com/long-mathematics/luce-tsetlin-permutations/blob/25b7e4490c32418b818b2356106228283986dd6c/papers/macroscopic-pd/poisson-dirichlet-luce.tex}{\texttt{macroscopic-pd}}.

\bibitem{LongSukhatme2026}
C. D. Long.
\newblock Sukhatme fixed points and the terminal Poisson process.
\newblock Preprint, 2026.
\newblock Repository manuscript:
\href{https://github.com/long-mathematics/luce-tsetlin-permutations/tree/main/papers/sukhatme-fixed-points}{\texttt{sukhatme-fixed-points}}.

\bibitem{Luce1959}
R. D. Luce.
\newblock \emph{Individual Choice Behavior: A Theoretical Analysis}.
\newblock Wiley, New York, 1959.

\bibitem{Pitman2006}
J. Pitman.
\newblock \emph{Combinatorial Stochastic Processes}.
\newblock Lecture Notes in Mathematics, vol. 1875.
\newblock Springer-Verlag, Berlin, 2006.

\bibitem{Plackett1975}
R. L. Plackett.
\newblock The analysis of permutations.
\newblock \emph{Journal of the Royal Statistical Society, Series C
(Applied Statistics)} 24 (1975), 193--202.
\newblock \href{https://doi.org/10.2307/2346567}
{doi:10.2307/2346567}.

\bibitem{Renyi1953}
A. R\'enyi.
\newblock On the theory of order statistics.
\newblock \emph{Acta Mathematica Academiae Scientiarum Hungaricae}
4 (1953), 191--231.
\newblock \href{https://doi.org/10.1007/BF02127580}
{doi:10.1007/BF02127580}.

\bibitem{SheppLloyd1966}
L. A. Shepp and S. P. Lloyd.
\newblock Ordered cycle lengths in a random permutation.
\newblock \emph{Transactions of the American Mathematical Society}
121 (1966), 340--357.

\bibitem{Sukhatme1937}
P. V. Sukhatme.
\newblock Tests of significance for samples of the $\chi^2$-population with
two degrees of freedom.
\newblock \emph{Annals of Eugenics} 8 (1937), no.~1, 52--56.
\newblock \href{https://doi.org/10.1111/j.1469-1809.1937.tb02159.x}
{doi:10.1111/j.1469-1809.1937.tb02159.x}.

\end{thebibliography}

\end{document}
